\documentclass[11pt,a4paper]{article}
\usepackage[margin=1.7cm]{geometry}
\usepackage{fontspec}
\IfFontExistsTF{Avenir Next}{
  \setmainfont{Avenir Next}
}{
  \IfFontExistsTF{Helvetica Neue}{
    \setmainfont{Helvetica Neue}
  }{
    \IfFontExistsTF{Source Sans 3}{
      \setmainfont{Source Sans 3}
    }{
      \setmainfont{Times New Roman}
    }
  }
}
\usepackage[hidelinks]{hyperref}
\setlength{\parindent}{0pt}
\setlength{\parskip}{0.35em}

\begin{document}

{\LARGE \textbf{Cristi Prodan}}\\
Telefon: +40 746 708 927 \\
Email: \href{mailto:prodan.cristian@gmail.com}{prodan.cristian@gmail.com} \\
GitHub: \href{https://github.com/christian}{https://github.com/christian} \\
LinkedIn: \href{https://www.linkedin.com/in/cristianprodan/}{https://www.linkedin.com/in/cristianprodan/} \\

\section*{Rezumat}
Senior Software Engineer specializat în platforme backend, sisteme distribuite și infrastructură cloud. Construiesc servicii de producție fiabile, le operez end-to-end și lucrez confortabil cu API-uri, securitate, CI/CD, observabilitate și livrare de produs.

\section*{Experiență}
\subsection*{Fondator, Comerz (Remote, România) \hfill Mar 2025 - prezent}
\begin{itemize}
  \item Am fondat și construit de la zero \href{https://comerz.ro}{https://comerz.ro}, o platformă e-commerce multitenant care îi ajută pe antreprenori să creeze și să administreze magazine online.
  \item Am proiectat și implementat MVP-ul end-to-end cu Rails și PostgreSQL: izolarea tenant-urilor, administrarea vitrinelor online, fluxuri pentru produse și catalog, plus deployment în producție.
  \item Am folosit fluxuri de dezvoltare asistate de AI cu Codex, Conductor și Linear pentru a accelera implementarea, iterația de produs și livrarea.
\end{itemize}

\subsection*{Senior Software Engineer, Zerofy (Remote, România) \hfill Ian 2024 - Feb 2025}
\begin{itemize}
  \item Am construit fluxuri backend pentru aplicația mobilă Zerofy folosind Google Cloud Functions, Firebase și Node.js pe GCP.
  \item Am implementat integrări serverless, event-driven, care conectau produsul cu furnizori externi și surse de date.
  \item Am reverse-engineered interfețe de vendor nedocumentate atunci când documentația oficială era incompletă sau indisponibilă.
  \item Am îmbunătățit mentenabilitatea prin refactorizări țintite, cleanup și limite mai clare în codul de integrare.
\end{itemize}

\subsection*{Senior Software Engineer, PTC Schweiz (Zürich, Elveția și Remote, România) \hfill Oct 2015 - Ian 2024}
\begin{itemize}
  \item Am construit, deținut și operat microservicii Vuforia Cloud în Scala/Akka pe AWS pentru produse de realitate augmentată folosite în manufacturing.
  \item Am colaborat cu ingineri de Machine Learning pentru integrarea modelelor în API-uri de producție și servicii containerizate, inclusiv un pipeline de date și orchestration workflow pentru modele ML construit de la zero.
  \item Am gestionat preocupări de sisteme distribuite: scalare, toleranță la erori, load balancing, securitatea serviciilor, deployment-uri de producție și răspuns la incidente on-call.
  \item Am operat servicii pe Kubernetes și am îmbunătățit observabilitatea cu Elasticsearch, Logstash și Kibana.
  \item Am construit fluxuri CI/CD și infrastructură folosind GitHub Actions, TeamCity, AWS CloudFormation și Terraform.
  \item Am livrat API-uri backend pentru clienți, inclusiv API-uri de căutare pentru Vivino și API-uri de producție pentru integrări de modele ML.
  \item Am crescut fiabilitatea prin teste end-to-end pentru REST APIs în Scala și refactorizări ale unor servicii mature.
  \item Am întărit securitatea prin integrarea Veracode și Snyk, aplicarea practicilor OWASP și implementarea Single Sign-On cu Amazon Cognito.
  \item Am construit tooling intern pentru deployment/productivitate și am condus inițiative care au îmbunătățit productivitatea dezvoltatorilor în mai multe echipe.
  \item Am scris documentație OpenAPI v3 și am implementat un serviciu proxy WebSocket pentru colaborare real-time în remote assistance.
  \item Am susținut interviuri tehnice și am contribuit la decizii de recrutare în engineering.
\end{itemize}

\subsection*{Senior Engineer, Qualcomm (Zürich, Elveția) \hfill Ian 2014 - Oct 2015}
\begin{itemize}
  \item Am extins servicii cloud AWS cu Scala și Play.
  \item Am automatizat operațiuni și tooling intern folosind Python și Bash.
  \item Am portat servicii către Docker și am menținut sisteme legacy Ruby on Rails.
  \item Am proiectat și construit un DSL intern în Ruby pentru testarea API-urilor între mai multe echipe.
  \item Am participat la scalarea infrastructurii de la câteva instanțe EC2 la sute de instanțe pe AWS (S3, SQS, DynamoDB, Cognito, Autoscaling).
\end{itemize}

\subsection*{Software Engineer, Kooaba (Zürich, Elveția) \hfill Feb 2011 - Ian 2014}
\begin{itemize}
  \item Am inițiat și condus migrarea de la un monolit la microservicii Scala, îmbunătățind scalabilitatea și reducând consumul de memorie.
  \item Am construit REST APIs și servicii backend cu Ruby on Rails.
  \item Am rescris și menținut o aplicație Android.
  \item Am optimizat query-uri și indexuri MySQL pe baze de date de producție cu peste 30 de milioane de rânduri.
\end{itemize}

\subsection*{Head of Technology, Alliants (Cluj-Napoca, România) \hfill Sep 2008 - Feb 2011}
\begin{itemize}
  \item Am angajat și coordonat o echipă de 3 ingineri care livra aplicații Ruby on Rails pentru proiecte ale clienților.
  \item Am coordonat livrarea on-site cu clienți din Londra și am transpus cerințele de business în planuri de implementare.
\end{itemize}

\section*{Competențe}
\begin{itemize}
  \item Limbaje: \textbf{Scala}, Ruby, Python, JavaScript, Bash
  \item Backend și API-uri: Akka/Pekko, Rails, Play, Node.js, REST APIs, OpenAPI, OAuth2, JWT, SSO
  \item Cloud și infrastructură: AWS, GCP, Docker, Kubernetes, CloudFormation, Terraform, Serverless
  \item Date și search: MySQL, PostgreSQL, DynamoDB, Elasticsearch, Logstash, Kibana
  \item Testare și tooling: GitHub Actions, TeamCity, Linux
\end{itemize}

\section*{Educație}
\begin{itemize}
  \item Master, Informatică --- Universitatea Babeș-Bolyai, Cluj-Napoca (2008-2010). Lucrare de disertație: Sisteme de recomandare folosind filtrare colaborativă. \href{https://github.com/christian/Rho-article}{https://github.com/christian/Rho-article}
  \item Licență, Informatică --- Universitatea Alexandru Ioan Cuza, Iași (2005-2008). Lucrare de licență: Sistem de recomandare mobil.
  \item Curs Machine Learning --- Andrew Ng (2011). Finalizat înainte ca platforma Coursera să fie lansată public.
\end{itemize}

\section*{Limbi}
Engleză (fluent), Română (nativ), Germană (nivel de bază)

\section*{Proiecte personale}
\begin{itemize}
  \item Virtual Private Organ: Proiect hardware/software pe termen lung pentru construirea unei orgi digitale personalizate. Include prelucrare lemn, electronică MIDI, programare Arduino, setup audio și documentație.
\end{itemize}

\end{document}
